\documentclass[10pt,handout]{beamer}

\input{EC2211_preamble.txt}
\title{Lecture Notes 2 \newline Economic Growth and the Solow Model}

\begin{document}

\maketitle


\begin{frame}
    \frametitle{Contents and literature}
    \begin{itemize}
        \item The production function
        \item Growth accounting
        \item The Solow model
    \end{itemize}\bigskip
    Literature: 
    \begin{tightitemize}
        \item Jones (2024) chapters 3-5
    \end{tightitemize}
\end{frame}


%%%%%%%%%%%%%%%%%%%%%%%%%%%%%%%%%%%%%%%%%%%%%%%%%%%
%%%%%%%%%%%%%%%%%%%%%%%%%%%%%%%%%%%%%%%%%%%%%%%%%%%
\section{Production}


\begin{frame}{The production function}
    \begin{block}{The Cobb Douglas production function}
        $Y = F(K,L) = A K^\alpha L^{1-\alpha}$
    \end{block} \bigskip
    
    \begin{table}
        \flushleft
        \begin{tabular}{rl}
            $Y$ & = output, GDP \\
            $F(.,.)$ & = production function \\
            $A$ & = total factor productivity (TFP)\\
            $K$ & = capital input \\
            $L$ & = labor input \\
            $\alpha$ & = capital share in production, $\alpha \in (0,1)$ \\
        \end{tabular}
    \end{table}
\end{frame}


\begin{frame}{Math repetition: derivatives}
Let $y=f(x)$ \par \medskip
The derivative of $y$ with respect to $x$ is written
\begin{equation*}
    \frac{\d{y}}{\d{x}}=\frac{\d{f(x)}}{\d{x}}=f'(x)
\end{equation*} \medskip
\end{frame}


\begin{frame}{Math repetition: derivatives}
    Let $y=f(x_1,x_2)$ \bigskip

    The derivatives of $y$ with respect to its arguments, $x_1$ and $x_2$, are written:
    \begin{equation*}
        \frac{\partial y}{\partial x_1}=\frac{\partial f(x_1,x_2)}{\partial x_1}=f'_{x_1}(x_1,x_2)
    \end{equation*}
    and
    \begin{equation*}
        \frac{\partial y}{\partial x_2}=\frac{\partial f(x_1,x_2)}{\partial x_2}=f'_{x_2}(x_1,x_2)
    \end{equation*}
\end{frame}


\begin{frame}{Math repetition: derivatives}
    Recall that if $y=f(x)=x^\alpha$, then $f'(x)=\alpha x^{\alpha-1}$ \bigskip

    Examples: 
    \begin{itemize}
        \item If $f(x)=x+0.5x^2-x^4 $, then $f'(x)=1+x-4x^3$
        \item If $f(x_1,x_2)=x_1^\alpha x_2^{1-\alpha}$, then $f'_{x_1}(x_1,x_2) = \alpha x_1^{\alpha-1} x_2^{1-\alpha}$
    \end{itemize}
\end{frame}


\begin{frame}{The production function}
\small{
    \begin{itemize}
        \item How much does production increase if we increase capital but hold labor constant?
        \item[] The marginal product of capital, $MPK$:
    \end{itemize}
    \[
    MPK = \frac{\partial Y}{\partial K} = \frac{\partial F(K,L)}{\partial K}= F_{K}
    \]
\bigskip

    \begin{itemize}
        \item How much does production increase if we increase labor but hold capital constant?
        \item[] The marginal product of labor, $MPL$:
    \end{itemize}
    \[
    MPL = \frac{\partial Y}{\partial L} = \frac{\partial F(K,L)}{\partial L}= F_{L}
    \]   
}
\end{frame}


\begin{frame}{Constant returns to scale}
    Consider again the Cobb-Douglas production function:
    \[
        Y=AK^\alpha L^{1-\alpha}
    \] 

    If all production factors are increased by a factor x, production also increases by factor x:
    \[
        F(xK, xL)=A(xK)^\alpha (xL)^{1-\alpha} = A x^{(\alpha+1-\alpha)}K^\alpha L^{1-\alpha}= xF(K,L)
    \] \pause

    \begin{itemize}
        \item If, for example, $F(K,L) = 10$ then $F(3K,3L) = 30$
        \item A function with this property is said to display \mfemph{constant returns to scale}
    \end{itemize}
\end{frame}


\begin{frame}{Diminishing marginal products}
\small{
    We still consider the Cobb-Douglas production function $Y=AK^\alpha L^{1-\alpha}$ \medskip

    The marginal products of capital and labor are
    \[
        MPK=\alpha AK^{\alpha-1} L^{1-\alpha} = \alpha A \left( \frac{L}{K} \right) ^{1-\alpha},
    \]
    and
    \[
        MPL=(1-\alpha) AK^\alpha L^{-\alpha} = (1-\alpha) A \left( \frac{K}{L} \right) ^\alpha
    \] \medskip

    \begin{itemize}
        \item If we increase $K$ but hold $L$ constant, production will increase but less and less as $K$ becomes larger
        \item If we increase $L$ but hold $K$ constant, production will increase but less and less as $L$ becomes larger
        \item The production function displays \mfemph{diminishing marginal products}
    \end{itemize}
}
\end{frame}


\begin{frame}{Diminishing marginal product of capital}
    \begin{figure}
        \includegraphics[scale=0.4]{LN02_Production_and_Solow/MACRO6_FIG04.01.jpg}
    \end{figure} \medskip
    \flushright \tiny Figure 4.1 in Jones (2024)
\end{frame}


\begin{frame}{Math repetition: profit maximization}
\begin{columns}
    \column{0.6\textwidth}
    \small{
    \begin{itemize}
        \item Consider a function $f(x)$
        \item Recall that $f'(x)$ is the slope of $f$ in point $x$
        \item If some $x^*$ is maximizes the function, then $f'(x^*)=0$
        \item This is the \mfemph{first-order condition} (f.o.c.)
        \item \footnotesize{(But note that the same holds for a minimum; the f.o.c. is a \emph{necessary} but not sufficient condition)}
    \end{itemize}
    }

    \column{0.4\textwidth}
    \begin{figure}
        \centering
        \includegraphics[width=1\linewidth]{LN02_Production_and_Solow/fx.png}
    \end{figure}
\end{columns}
\end{frame}


\begin{frame}{Profit maximization under Cobb-Douglas}
    Competitive firms choose capital and labor to maximize profits:
    \begin{equation*}
        \max_{K,L} \Pi(K,L) =\underbrace{AK^\alpha L^{1-\alpha}}_{F(K,L)}-rK-wL
    \end{equation*}

    where $\Pi$ denotes the profit function, $w$ the wage and $r$ the rental rate of capital
    \bigskip

    First-order conditions for profit maximization:
    \begin{equation} \label{focK}
        \frac{\partial \Pi}{\partial K} = \underbrace{\alpha AK^{\alpha-1} L^{1-\alpha}}_{MPK}-r=0
    \end{equation}
    \begin{equation} \label{focL}
        \frac{\partial \Pi}{\partial L} = \underbrace{(1-\alpha) AK^{\alpha} L^{-\alpha}}_{MPL}-w=0
    \end{equation}
\end{frame}


\begin{frame}{The demand for capital and labor}
    Re-arranging \eqref{focK}:
    \begin{equation} \label{MPK}
        MPK=\alpha A K^\alpha L^{(1-\alpha)} K^{-1}=\alpha \frac{Y}{K}=r
    \end{equation}
    
    Re-arranging \eqref{focL}:
    \begin{equation} \label{MPL}
        MPL= (1-\alpha) A K^\alpha L^{(1-\alpha)} L^{-1} = (1-\alpha)\frac{Y}{L}=w
    \end{equation}
\end{frame}


\begin{frame}{Equilibrium in the capital and labor markets}
    \begin{figure}
        \includegraphics[scale=0.4]{MACRO6_FIG04.02.jpg}
    \end{figure} \medskip
    \flushright \tiny Figure 4.2 in Jones (2024)
\end{frame}


\begin{frame}{}
    Note that \eqref{MPK} and \eqref{MPL} imply:
    \begin{equation*}
        \alpha=\frac{rK}{Y}
    \end{equation*}

    and

    \begin{equation*}
        1-\alpha=\frac{wL}{Y}
    \end{equation*}\bigskip

    Jones (2024) assumes that $\alpha=1/3$. Why?
\end{frame}



\begin{frame}{Labor's share of GDP ($1-\alpha$)}
    \begin{figure}
        \includegraphics[scale=0.4]{MACRO6_FIG02.03.jpg}
    \end{figure}
    \flushright \tiny Figure 2.3 in Jones
\end{frame}


\begin{frame}{Labor's share of GDP}
\small{
\begin{itemize}
    \item Macroeconomists often say that the labor's share of GDP in the United States has been roughly constant at around $2/3$ for the last century \pause
    
    \item Data for the U.S. and many other countries however indicate that the labor's share has declined somewhat since the 1980s (see previous slide and \href{https://www.clevelandfed.org/-/media/project/clevelandfedtenant/clevelandfedsite/publications/economic-commentary/2018/ec-201808-evolution-of-the-labor-share-across-developed-countries/ec_201808.pdf}
    {this note})
    \begin{tightitemize}
        \item The reasons behind this fall are not obvious
        \item Separating compensation to labor from compensation to capital owners is not always trivial
    \end{tightitemize} \pause
    
    \item Sweden may be an exception (has no clear trend)
    \begin{tightitemize}
        \item See next slide and \href{https://ekonomistas.se/2024/01/17/konjunkturinstitutet-replik-om-vinstandelens-utveckling/}{this blog post}
        \item[] \footnotesize{(The slide and the blog post are in Swedish. Some translations: "vinstandelar" = "profit shares", "hela ekonomin" = "total economy", "näringslivet" = "private sector")}
    \end{tightitemize}
\end{itemize}
}
\end{frame}


\begin{frame}{Capital's share of GDP, Sweden ($\alpha$)}
\begin{figure}
  \includegraphics[scale=0.24]{LN02_Production_and_Solow/profitshare_se.png}
\end{figure}
\flushright{\tiny Source: \href{https://ekonomistas.se/2024/01/17/konjunkturinstitutet-replik-om-vinstandelens-utveckling/}{ekonomistas.se}}
\end{frame}




%%%%%%%%%%%%%%%%%%%%%%%%%%%%%%%%%%%%%%%%%%%%%%%%%%%
%%%%%%%%%%%%%%%%%%%%%%%%%%%%%%%%%%%%%%%%%%%%%%%%%%%
\section{Growth accounting}


\begin{frame}{Growth accounting}
    Can we use the production function to "explain" why some countries are so much richer than others? \pause \medskip

    \begin{itemize}
        \item Consider our Cobb-Douglas production function:
        \[
            Y_t = A_t K_t^\alpha L_t^{1-\alpha}
        \]

        \item Apply the rules for growth rates:
        \[
            g_{Y} = g_{A} + \alpha g_{K} + (1-\alpha) g_{L}
        \]
    
        \item This shows that the growth rate of output is the sum of TFP growth and (weighted) growth of the production factors \smallskip

        \pause\item We can observe or estimate $Y_t$, $K_t$, $L_t$ and $\alpha$ in the data \smallskip
    
        \item Growth in TFP is then calculated as the residual
        \[
            g_{A} = g_{Y} - \alpha g_{K} - (1-\alpha) g_{L}
        \]
    \end{itemize}
\end{frame}


\begin{frame}{Income per capita}

    Cross-country comparisons require expressions in per-capita terms\bigskip

    Recall the production function $Y=A K^\alpha L^{1-\alpha}$ and divide by $L$:
    \begin{equation*}
        y \equiv \frac{Y}{L}=\frac{AK^\alpha L^{1-\alpha}}{L}=A K^\alpha L^{-\alpha}=Ak^\alpha,
    \end{equation*}
    where $y=Y/L$ and $k \equiv K/L$ denote GDP and capital per person ("per capita"), respectively
\end{frame}

\begin{frame}{Growth accounting} 
    In per capita terms, production is
        \[ y_t = A_t k_t^\alpha \]
    so:
    \[
        g_{y} = g_{A} + \alpha g_{k}
    \] \medskip

    If we compare country $i$ to the United States:
    \[
        \frac{y_i}{y_{US}} = \frac{A_i k_i^\alpha}{A_{US} k_{US}^\alpha} = \left( \frac{A_i}{A_{US}} \right) \left( \frac{k_i}{k_{US}} \right) ^\alpha
    \]
\end{frame}





\begin{frame}{Predicted per capita GDP}
\begin{itemize}
    \item Suppose $y=Ak^{1/3}$ and that $A$ is the same as in the United States in all countries
    \bigskip
    \item Actual differences in production are much larger than predicted by differences in capital input {\tiny (see next two slides!)}
    \bigskip
    \item If this is the production function, much of differences in GDP per capita must be attributed to differences in TFP
\end{itemize}
\end{frame}


\begin{frame}{Predicted per capita GDP relative to the U.S. if $y=k^{1/3}$}
    \begin{figure}
        \centering
        \includegraphics[width=0.79\linewidth]{MACRO6_Table04.03_box.jpg}
    \end{figure} \pause

    \footnotesize{
    \begin{tightitemize}
        \item Consider, for example, China: the table shows that $k^{China}/k^{USA} = 0.275$
        \item If $y=k^{1/3}$, then $y_{China}/y_{US} = (k_{China}/k_{US})^{1/3} = 0.275^{1/3} = 0.650 $
        \item Actual GDP per capita in China is just 22.6 percent of GDP per capita in the United States
    \end{tightitemize}
    }
\end{frame}


\begin{frame}{Predicted per capita GDP}
\begin{figure}
  \includegraphics[scale=0.30]{Fig_implied_gdp.png}
  \caption{\tiny{This graph shows similar information as the table on the previous page but for more countries. Horizontal axis: Actual GDP per capita in 2019 (2017 USD) from PWT 10.01. Vertical axis: Predicted GDP per capita according to $Ak^{1/3}$ where $A$ is based on the U.S. value and $k$ calculated from PWT with the same method as Jones. Ireland and Qatar were dropped from sample.}}
\end{figure}
\end{frame}


\begin{frame}{Much more differences in GDP per capita than in capital per capita}
    \begin{figure}
        \centering
        \includegraphics[width=0.5\linewidth]{WorldBank2024.png}
    \end{figure}
    \flushright \tiny{Figure O.2 in \href{https://bit.ly/WDR2024_FullReport_ENG}{World Bank (2024)}}
\end{frame}


\begin{frame}{If capital is the only production factor, TFP in most other countries must be much lower than in the U.S. to match the data}
    \begin{figure}
        \centering
        \includegraphics[width=0.85\linewidth]{MACRO6_Table04.04_box.jpg}
    \end{figure}
    \flushright \tiny{Table 4.4 in Jones}
\end{frame}


\begin{frame}{Why does TFP differ across countries?}
\begin{itemize}
    \item The accounting exercise in Jones (2024) suggests that $A$ is twice as important as $k$ in explaining cross-country differences in $y$ \bigskip
    
    \item Why does TFP differ across countries? \smallskip
    \begin{tightitemize}
        \item (Labor input)\smallskip
        \item Human capital\smallskip
        \item Technology\smallskip
        \item Institutions\smallskip
        \item Misallocation
    \end{tightitemize}
\end{itemize}
\end{frame}


\begin{frame}{Another indication of the importance of TFP}
\begin{itemize}
    \item Suppose that all countries have identical production functions $Y=AK^\alpha L^{1-\alpha}$ and identical $A$ \smallskip

    \item Suppose that physical capital can move freely between countries\smallskip

    \item Then capital should flow from rich to poor countries because the MPK would then be much higher in poor countries {\footnotesize (\href{https://www.jstor.org/stable/2006549}{Robert Lucas, 1990, "Why doesn't capital flow from rich to poor countries"})} \smallskip

    \item The current account balance shows capital flows: $CA=S-I$
    \begin{tightitemize}
        \item A negative current account balance means that capital flows to the country (investment is larger than saving)
        \item There is some tendency for capital to flow to poor countries
        \item But not at all to the extent this "model" predicts
    \end{tightitemize}
\end{itemize}
\end{frame}


\begin{frame}{Some tendency for capital to flow from rich to poor countries in 2010s}
\begin{figure}
  \includegraphics[scale=0.15]{Fig_capitalflows_2010_2019.png}
\end{figure}
\tiny Averages for 2010-2019. Current account balance in percent of GDP. GDP per capita relative to USA. Source: IMF WEO.
\end{frame}


\begin{frame}{Modest capital flows in 1980s}
\begin{figure}
  \includegraphics[scale=0.15]{Fig_capitalflows_1980_1989.png}
\end{figure}
\tiny Averages for 1980-1989. Current account balance in percent of GDP. GDP per capita relative to USA. Source: IMF WEO.
\end{frame}




%%%%%%%%%%%%%%%%%%%%%%%%%%%%%%%%%%%%%%%%%%%%%%%%%%%
%%%%%%%%%%%%%%%%%%%%%%%%%%%%%%%%%%%%%%%%%%%%%%%%%%%
\section{The Solow model}


\begin{frame}{The Solow model}
\begin{itemize}
    \item The production function is the key ingredient\medskip

    \item Does capital accumulation generate growth?\medskip

    \item A key model for\smallskip
    \begin{tightitemize}
        \item ... understanding economic growth and its limitations\smallskip
        \item ... but also a building stone for much other macroeconomic analysis
    \end{tightitemize}
\end{itemize}
\end{frame}


\begin{frame}{The Solow model}
\begin{columns}
    \column{0.6\textwidth}
    \begin{itemize}
    \item Solow (1956), "A Contribution to the Theory of Economic Growth"\medskip

    \item \href{https://www.nobelprize.org/prizes/economic-sciences/1987/solow/facts/}{Nobel prize 1987}: \small{"The study of the factors which permit production growth and increased welfare has been a central feature in economic research for many years. Robert M. Solow’s prize recognizes his exceptional contributions in this area."}
    \end{itemize}

    \column{0.4\textwidth}
    \begin{figure}
        \centering
        \includegraphics[width=0.9\linewidth]{solow-13390-portrait-mini-2x.jpg}
    \end{figure}
    \centering \small{Robert Solow 1924-2023}
\end{columns}
\end{frame}


\begin{frame}{Notation}
\begin{itemize}
    \item[] $Y_t$: output at time $t$
    \item[] $A$: total factor productivity (TFP)
    \item[] $K_t$: the capital stock at time $t$
    \item[] $L_t$: labor input at time $t$
    \item[] $I_t$: investment at time $t$
    \item[] $C_t$: consumption at time $t$ 
    \item[] $K_0$: initial stock of capital
    \item[] $\delta$: the depreciation rate of capital\footnote{\tiny{Jones uses $\overline{d}$ to denote the depreciation rate, but $\delta$ is a much more common notation }}
    \item[] $s$: savings rate
\end{itemize}
\end{frame}


\begin{frame}{Production}
    Production is a function of the two production factors, capital and labor:
    \begin{equation} \label{eq1}
        Y_t=F(K_t, L_t)
    \end{equation}\smallskip

    We assume that the production function exhibits \emph{constant returns to scale} and \emph{diminishing marginal products of capital and labor}.
    \smallskip
    Closed-economy resource constraint: 
    \begin{equation} \label{eq2}
        C_t+I_t=Y_t
    \end{equation}
\end{frame}


\begin{frame}{Capital accumulation}
    Capital in period $t+1$ is determined by:
    \begin{equation} \label{eq3}
        K_{t+1}=K_t+I_t-\delta K_t
    \end{equation} \smallskip

    Equation \eqref{eq3} can be re-arranged to read:
    \begin{equation} \label{eq4}
        \Delta K_{t+1}= \underbrace{I_t - \delta K_t}_{\text{net investment}}
    \end{equation}
    where $\Delta K_{t+1} \equiv K_{t+1}-K_t$\medskip

    In each period $t>0$, $K_t$ results from previous investment and depreciation\smallskip

    To get the model started, an initial value for $K_0$ is needed
\end{frame}


\begin{frame}{Consumption and investment}
    Assume that a constant fraction of income is saved (and invested):
    \begin{equation} \label{eq5}
        I_t=sY_t
    \end{equation}\smallskip

    This implies that consumption is:
    \begin{equation} \label{eq6}
        C_t=(1-s)Y_t
    \end{equation}
\end{frame}


\begin{frame}{Solving the model}
    Assume for now that there is no population growth so that labor input is also constant over time, $L_t=L$, and we also assume that TFP is constant over time, $A_t=A$  \bigskip

    Equations \eqref{eq4}, \eqref{eq5} and \eqref{eq1} imply:
    \begin{equation} \label{eq7}
        \Delta K_{t+1}= sF(K_t,L) - \delta K_t
    \end{equation}
\end{frame}


\begin{frame}{Cobb-Douglas production function}
    Assume that the production function is Cobb-Douglas: 
    \begin{equation} 
        Y_t=A K_t^\alpha L^{1-\alpha}
    \end{equation} \smallskip

    \begin{itemize}
    \item Recall that the marginal product of capital is then diminishing:
    \begin{tightitemize}
        \item We have assumed that labor input is fixed
        \item Plotted against capital, $Y$ and $sY$ will have the concave shape we saw previously
    \end{tightitemize}
    \item Depreciation, $\delta K$, is however \emph{linear}
    \item The change in the capital stock from one period to the next is investment minus depreciation, i.e. the difference between $sY$ and $\delta K$
    \end{itemize}
\end{frame}


\begin{frame}{The Solow Diagram}
    \begin{figure}
        \centering
        \includegraphics[width=0.85\linewidth]{SolowDiagram1.png}
    \end{figure}
\end{frame}


\begin{frame}{Steady state}
    \begin{block}{Steady state}
        In a steady state, all variables are constant
    \end{block} \bigskip \pause

    If capital is constant, then $K_{t+1}=K_t$, i.e. $\Delta K_{t+1}=0$ \bigskip

    Equation \eqref{eq7} implies that the steady-state level of capital $K^{*}$ must satisfy: 
    \begin{equation} \label{eq8}
        sY^* = sF(K^{*},L)=\delta K^{*}
    \end{equation}
    where $Y^*$ is steady-state output
\end{frame}


\begin{frame}{Steady state under Cobb-Douglas}
    We assume that the production function is Cobb-Douglas, 
    \begin{equation} \label{eq9}
        Y_t=A K_t^\alpha L^{1-\alpha}
    \end{equation} \smallskip

    Equation \eqref{eq8}, $sY^*=\delta K^*$, then implies that
    \begin{equation*}
        sA (K^{*})^\alpha L^{1-\alpha} = \delta K^{*}
    \end{equation*}
\end{frame}


\begin{frame}{Steady state under Cobb-Douglas}
    Solving for $K^{*}$: 
    \begin{equation} \label{eq10}
        K^{*}=\bigg( \frac{sA}{\delta} \bigg)^{\frac{1}{1-\alpha}} L
    \end{equation} \smallskip

    Plugging \eqref{eq10} into \eqref{eq9} yields steady-state output:
    \begin{equation} \label{eq11}
        Y^{*} = A^{\frac{1}{1-\alpha}} \bigg( \frac{s}{\delta} \bigg)^{\frac{\alpha}{1-\alpha}} L
    \end{equation}
\end{frame}


\begin{frame}{Steady-state capital and output per worker}
    Dividing \eqref{eq10} by $L$ we obtain steady-state capital per worker: 
    \begin{equation} \label{eq12}
        k^{*} \equiv \frac{K^{*}}{L} = \bigg( \frac{sA}{\delta} \bigg)^{\frac{1}{1-\alpha}}
    \end{equation} \smallskip

    Dividing \eqref{eq11} by $L$ we obtain steady-state output per worker: 
    \begin{equation} \label{eq13}
        y^{*} \equiv \frac{Y^{*}}{L} = A^{\frac{1}{1-\alpha}} \bigg( \frac{s}{\delta} \bigg)^{\frac{\alpha}{1-\alpha}}
    \end{equation}
\end{frame}


\begin{frame}{The Solow Diagram}
    \begin{figure}
        \includegraphics[scale=0.4]{MACRO6_FIG05.01.jpg}
    \end{figure} \medskip
    \flushright \tiny Figure 5.1 in Jones (2024)
\end{frame}


\begin{frame}{The Solow Diagram with output}
    \begin{figure}
        \includegraphics[scale=0.4]{MACRO6_FIG05.02.jpg}
    \end{figure} \medskip
    \flushright \tiny Figure 5.2 in Jones (2024).
\end{frame}


\begin{frame}{Exercise: Where is the steady state? (Poverty trap?)}
    \begin{figure}
        \includegraphics[scale=0.48]{PovertyTrap2.png}
    \end{figure}
    \tiny Note that this production function is non-standard and not Cobb Douglas -- we use it here just to illustrate how investment and depreciation determine the dynamics of the capital stock
\end{frame}


\begin{frame}{Some model implications}
\begin{itemize}
    \item More capital accumulation (higher $s$) leads to higher income but not to sustained growth\bigskip

    \item Why does the economy reach a steady state?
    \begin{tightitemize}
        \item Diminishing returns to capital, but depreciation is proportional to capital
        \item In the steady state, investment is just offset by capital depreciation
    \end{tightitemize} \bigskip

    \item The capital-output ratio in steady state is proportional to the savings rate \bigskip

    \item The savings rate can be inefficiently high ("Golden Rule", see exercise)
\end{itemize}
\end{frame}


\begin{frame}{An increase in the investment rate}
\pause
    \begin{figure}
        \includegraphics[scale=0.4]{MACRO6_FIG05.05.jpg}
    \end{figure} \medskip
    \flushright \tiny Figure 5.5 in Jones (2024)
\end{frame}


\begin{frame}{Capital-output ratio and savings}
    Recall that equation \eqref{eq8} predicts $sY^{*}=\delta K^{*}$, so that:
    \begin{equation} \label{eq14}
        \frac{K^{*}}{Y^{*}} = \frac{s}{\delta}
    \end{equation}

    \bigskip

    Idea: plot capital-income ratios against investment rates
\end{frame}


\begin{frame}{"Explaining" capital in the Solow model}
    \begin{figure}
        \includegraphics[scale=0.3]{MACRO6_FIG05.03.jpg}
    \end{figure}
    \flushright  \tiny {Figure 5.3 in Jones (2024)}
\end{frame}


\begin{frame}{Capital-output ratio and savings}
    \begin{itemize}
        \item Figure 5.3 looks like very strong support for the model in the data\medskip
        \item But how do we measure capital stocks?\medskip
        \pause
        \item Typically measured indirectly from investment rates\medskip
        \item Jones assumes that $\delta=0.05$ and calculates
        \begin{tightitemize}
            \item $K_{1960}=K^*$
            \item $K_{1961}=(1-\delta)K_{1960}+I_{1960}$, etc
        \end{tightitemize}
    \end{itemize}
\end{frame}


\begin{frame}{Transition dynamics}
\begin{itemize}
    \item No long-run growth in this version of the Solow model: the economy settles down at $K^{*}$ and $Y^{*}$

    \item But countries can grow temporarily during transitions to steady states
    \begin{tightitemize}
        \item If $K_t<K^{*}$ the growth rate will be positive
        \item If $K_t>K^{*}$ the growth rate will be negative
        \item The further from the steady state, the larger the (absolute value of) growth rates
    \end{tightitemize}

    \item The growth that we observe may be due to shocks that cause deviations from steady states, or to shifts in steady states
\end{itemize}
\end{frame}


\begin{frame}{Transitional dynamics, convergence}
\begin{itemize}
    \item Suppose that countries have similar production functions \& similar TFP \medskip
    \begin{tightitemize}
        \item They then have similar steady states
        \item But maybe some (the "poor") just have little capital
        \item Poor countries are then far from steady state, should grow rapidly
    \end{tightitemize} \medskip
    \item What does the data say?
\end{itemize}
\end{frame}


\begin{frame}{Growth rates in the OECD, 1960-2019}
    \begin{figure}
        \includegraphics[scale=0.35]{MACRO6_FIG05.08.jpg}
    \end{figure}
    \flushright \tiny Figure 5.8 in Jones (2024)
\end{frame}


\begin{frame}{Growth rates around the world, 1960-2019}
    \begin{figure}
    \includegraphics[scale=0.35]{MACRO6_FIG05.09.jpg}
    \end{figure}
    \flushright \tiny Figure 5.9 in Jones (2024)
\end{frame}


\begin{frame}{Convergence}
\begin{itemize}
    \item Clear evidence of convergence among OECD countries since 1960\medskip
    \item Less clear evidence of convergence in larger sample of countries\medskip
    \item Possible interpretation: 
    \begin{tightitemize}
        \item OECD countries are similar and have similar steady-states
        \item But this is not the case for other economies; they may be close to their steady-states
    \end{tightitemize}\medskip
    \item Convergence vs. conditional convergence
\end{itemize}
\end{frame}


\begin{frame}{South Korea and the Philippines 1960-2019}
\footnotesize{
\begin{table}
        \centering
        \begin{tabular}{L{2cm}C{2cm}C{2cm}C{2.5cm}}
         &GDP per capita relative to USA 1960&GDP per capita relative to USA 2019&Increase in GDP per capita 1960-2019 (percent)\\
         \hline
         USA& 1.00 &1.00&228\\
         Philippines& 0.11&0.13&317\\
         Korea&0.07&0.67&3261\\
    \end{tabular}
\end{table}
} \bigskip \pause

\normalsize {Can the Solow model explain why growth was so much more rapid in Korea than in the Philippines?}
\end{frame}


\begin{frame}{Case study: South Korea and the Philippines 1960-2019}
    \begin{figure}
        \includegraphics[scale=0.42]{MACRO6_FIG05.10_box.jpg}
    \end{figure}
    \flushright \tiny Figure 5.10 in Jones (2024)
\end{frame}


\begin{frame}{Summing up}
\begin{itemize}
    \item We have analyzed the basic Solow model with constant $L$ and $A$\medskip
    \item The model shows that capital accumulation:
    \begin{tightitemize}
        \item can generate transitional growth
        \item cannot generate sustained growth
        \item can be inefficiently high (exercise)
    \end{tightitemize}\medskip
    \item Taking the model to the data, much cross-country variation remains unexplained\medskip
    \item Next time: population growth, productivity, human capital, natural resources, ... 
\end{itemize}
\end{frame}

\end{document}


